% 2026-11-23 Michelman 419--424 頁（PDF 頁 23--28）與註 59--72，主對話以頁圖核對 419、421 頁之註號；
% 422--424 頁註 68--72 依文字層。Dworkin、Rawls、Habermas、Post《Constitutional Domains》、Raz 均無本機藏本，一律轉引 Michelman 之註。
\chapter{憲政民主之悖論與民主地追求民主正義}
\label{ch:40_paradox}

\begin{qaquestion}
Michelman\ix{persons}{michelman@Michelman, Frank I.} 在 419 頁原文第 V 部寫道：「No procedure can decide the question of that procedure's own fitness to decide either that or any other kind of question」〔沒有任何程序能夠決定該程序本身是否適於決定該問題或任何其他問題〕。
\begin{enumerate}
\tightlist
\item 這句話依據什麼理由成立？由它如何推出把根本法的決定劃出民主程序之外的衝動出於邏輯（420 頁）？
\item 原文第 VI 部（420--421 頁）以什麼為起點，區分了哪兩個問題？合理的詮釋多元論與 Rawls 的合理多元論有何差別（421--422 頁）？
\item 423 頁提出的正當性之候選屬性是什麼？Michelman 為何說這項回應不完整？
\item 「democratic pursuit of democratic justice」在 424 頁指怎樣的計畫？
\end{enumerate}
\end{qaquestion}

\section{程序不能決定自身之適格}
\label{sec:40_paradox-1}

原文第 V 部以一個設問起頭：關心民主的人必須「mean something by it」〔對民主有所指〕，他對民主的承諾，就是使本國的「laws of lawmaking」〔造法之法〕\ix{terms}{02@\zhlidxcat{憲政主義}!造法之法}符合他所謂民主的要求，因此「How can you possibly---how can you conceivably---leave democracy to decide the question of what democracy means in the first place?」〔怎麼可能，怎麼設想得出，讓民主來決定民主一開始是什麼意思這個問題？〕\footcite[419]{michelman1998}

Michelman 說，人們總可以敘述現行憲法如何取得其規範內容，也總可以主張那段歷史是民主的，但無論那段歷史多麼精緻地民主，都不能拿來作為肯認現行憲法、或要求他人尊重它的根據；要這麼做，須以「it is right to enforce on a country whatever constitutional law may have issued from that particular sort of assertedly democratic history」〔把出自那種自稱民主之歷史的任何憲法強制施行於一國是正當的〕為前提，而這個前提本身若再以它出自民主歷史這一事實為據，就陷入無窮後退（bad infinite regress）\ix{terms}{02@\zhlidxcat{憲政主義}!無窮後退}。註 59 引 Dworkin\ix{persons}{dworkin@Dworkin, Ronald} 為佐：民主「cannot prescribe the procedures for testing whether the conditions for the procedures it does prescribe are met」〔不能規定檢驗其所規定之程序的條件是否具備的程序〕。\footcite[419 \& n.59（轉引 Dworkin, Freedom's Law 33 (1996)）]{michelman1998}

Michelman 稱民主是「a demanding normative idea, an idea with content, however uncertain or disputable that content may be」〔一個要求甚高的規範理念，一個有內容的理念，不論那內容多麼不確定或有爭議〕。人們或許都同意民主意指由許多人共同行動的聯合決定程序，但它「no less essentially」〔同樣根本地〕也意指程序各方之間一種經社會構成的關係：參與者須在彼此之間、並相對於政治集體及其權力，以平等、獨立、自由與安全的適當關係進入程序，否則不稱其為民主程序。Michelman 把這層意思歸為 Dworkin 的 darker-hued message〔色調較暗的訊息〕。其結論是，沒有民主可以不先設定一份「to all intents and purposes a bill of rights」〔就所有實際目的而言等於權利清單〕，其操作內容須與這些關係相應。\footcite[419]{michelman1998}註 60 說明這項設定不涉及適用或「enforcing」〔執行〕該權利清單的機關安排，所以由法院執行並非這句話的內容。\footcite[419 n.60]{michelman1998}

權利清單的內容依某種大致確定的理解而定，420 頁接著說，那內容「can never be allowed to vary with the outcome of any vote or other collective-decision procedure」〔絕不容許隨任何投票或其他集體決定程序的結果而變動〕；註 61 以 Rawls 所稱純粹程序正義（pure procedural justice）作對照。\footcite[420 \& n.61]{michelman1998}同頁又說，不論在任何人眼中屬理想的民主條件下舉行的投票，也不保證該程序依這項理解所要求的方式決定特定的根本法問題。由此「evidently follows」〔顯然推得〕，關心民主的人須獨立於任何號稱民主的程序背書，自行負起決定這些問題的責任，或把責任交給法官，「Or to the philosophers」〔或交給哲學家〕。\footcite[420]{michelman1998}

於是 Michelman 說，把根本法的決定自民主的程序範圍劃出去的那股難以抑制的衝動（impulse to hive off，劃出）\ix{terms}{02@\zhlidxcat{憲政主義}!劃出}，出於必然，與民主做什麼才審慎或可欲的反省無涉：「it springs from logic」〔它來自邏輯〕。這套邏輯「entirely bars democracy from a decision-space」〔完全禁止民主進入一個決定空間〕，即決定一國最基本的法律、其 laws of lawmaking 的內容之處。\footcite[420]{michelman1998}原文第 V 部標題所稱憲政民主之悖論\ix{terms}{02@\zhlidxcat{憲政主義}!憲政民主之悖論}，依這段論證表現為：民主須以權利清單為前提，而權利清單的內容不能由民主決定。

\section{力量問題與兩個問題的區分}
\label{sec:40_paradox-2}

原文第 VI 部的標題帶問號：「The Democratic Pursuit of Democratic Justice?」。Michelman 說他要先簡短考慮，如何能在理論上重新證成民主程序參與決定一國最基本的法律，證成的不只是邏輯上的可能，還有道德上的必要。起點取自他所認的自由主義政治思想的源頭，即力量問題（the problem of power and force）：任何既立的政治秩序，遲早會對其權威範圍內的個人施以武力的威脅。自由主義思想視這些個人原初即自由而平等，因此推論出，這種武力的體制只有在受其拘束的每個人都有自己的理由同意其基本條款時才是正當的，其中包括它的 laws of lawmaking。\footcite[420]{michelman1998}

註 62 引 Rawls\ix{persons}{rawls@Rawls, John}《Political Liberalism》217 頁：政治權力的行使僅在依憲法行之、而憲法的要點可合理期待所有公民依其所接受的合理且理性之原則與理想予以認可時，才屬可證成。\footcite[420 n.62]{michelman1998}Michelman 據此說，這類法律要在自由主義的觀點下正當，就須回應每個人的理由（縱使有人未自覺這些理由），那些理由客觀上符合每個人的利益，或合於每個人作為合理的人、在適當的規範性論述中所會同意者；註 63 參見 Habermas\ix{persons}{habermas@Habermas, Jürgen}《Between Facts and Norms》。\footcite[420--421 \& n.63]{michelman1998}依此構想的正義為憲法設定一項客觀標準，現實中有的憲法內容符合，有的肯定不符。\footcite[421]{michelman1998}

哪一種憲法內容符合該標準，既是「a question not finally decidable by the procedure of democratic self-government」〔民主的自治程序無法終局決定的問題〕，Michelman 便說，這個關於真正正義的客觀問題，也不可能是自由主義公民最首要的政治道德問題；首要的問題是：我們以日常的合作行為支持現行政權，這種支持在道德上能否證成。這兩個問題相關，但「for liberals, the two questions cannot be identical」〔對自由主義者而言，兩者不可能同一〕。\footcite[421]{michelman1998}兩者不能重疊的理由，是他稱為合理的詮釋多元的事實。

\section{合理的詮釋多元論}
\label{sec:40_paradox-3}

Michelman 說，他依 Rawls 的用語風格，稱之為「the fact of reasonable interpretive pluralism」〔合理的詮釋多元的事實〕\ix{terms}{04@\zhlidxcat{正當性}!合理的詮釋多元}：一國之內大致自由的居民，對正義所要求的根本法解釋，無可避免地存在無法解決的不確定性與無法修補的合理歧見。\footcite[421 \& n.64]{michelman1998}註 64 指示參照其前文原文第 III 部 C 節與第 IV 部 B 節。他同時聲明並未嚴格依循 Rawls，兩人所說的多元，所涉層次不同，表 \ref{tab:40_paradox-pluralism} 依 Michelman 421--422 頁及註 65--69 並列。

\begin{longtable}{@{}L{0.14\linewidth}L{0.38\linewidth}L{0.38\linewidth}@{}}
\caption{Rawls 之合理多元論與 Michelman 之合理的詮釋多元論。}\label{tab:40_paradox-pluralism}\\
\toprule
項目 & Rawls 之合理多元 & Michelman 之合理的詮釋多元 \\
\midrule
\endfirsthead
\toprule
項目 & Rawls 之合理多元 & Michelman 之合理的詮釋多元 \\
\midrule
\endhead
多元之對象 & 整全性倫理觀點（comprehensive views）。註 65 引 Rawls 13 頁：對人生中何者有價值、人格理想、友誼與家庭及結社關係等的觀點 & 對自由主義憲法原則所作、道德上關鍵之解釋。這些原則是 Rawls 所稱倫理多元的各參與者皆有理由同意者，並已把其倫理多元之事實計入 \\
在理論中之地位 & 屬「general facts」〔一般事實〕，亦稱「circumstances of justice」〔正義之環境〕，影響人們對於涉及憲法正義內容之公共實踐規範，有無同意之理由（註 66 引 Rawls 36--37、66 頁：例如使人有理由接受確立之寬容規範） & 此種多元連同判斷之負擔，「cuts deeper」〔切入更深〕（421 頁） \\
判斷之負擔\ix{terms}{04@\zhlidxcat{正當性}!判斷之負擔} & 使通情達理之人對正義仍無法消除歧見的原因。註 67 引 Rawls 54--55 頁：這些人誠實觀察與陳述、有說服力地論辯，並共有「desire to honor fair terms of cooperation」〔遵守公平合作條件之願望〕；56--57 頁舉各人全部經驗必然不同，故判斷於相當複雜之案件上分歧 & 引 Rawls 之說，未另作界定 \\
所涉之真理 & Michelman 未就 Rawls 作對照 & 使憲法及其解釋對正義有何要求，任何權威的真理都無法經公共理性證明（422 頁）。真理在哲學上仍可得，也仍在理由論證的範圍之內；它在政治上不可得，因為各方既不能就真理為何達成一致，也無法把一切與己相反的立場視為不合理 \\
\bottomrule
\end{longtable}

註 69 載明，這種把合理多元自整全性倫理觀點降到政治正義原則之解釋層次的作法，對 Rawls 的「political liberal」〔政治自由主義〕回應有顛覆性的涵義，Michelman 說這點使他不安（makes me nervous）。\footcite[422 n.69]{michelman1998}註 68 另引 Raz\ix{persons}{raz@Raz, Joseph}，Raz 曾類似地懷疑，是否有理由認為人在美好生活的性質上，比在人人皆同意應影響政治行動的那類道德考量（如生命權、表現自由、宗教崇拜自由）上更可能出錯。\footcite[422 n.68]{michelman1998}

合理的詮釋多元在真正正義的問題與公民對政治武力應當怎麼做的問題之間，「opens a gap」〔開出一道缺口〕。更糟的是，自由主義者肯認合理的詮釋多元，等於讓自己面對一種可能：對政治武力應如何處置，根本沒有答案，凡是對它所做的都不可能正當，「all there can be is facts of power」〔所能有的只是權力的事實〕。Michelman 說這就是 Rawls 所稱「the problem of political liberalism」〔政治自由主義的問題〕。\footcite[422 \& n.69]{michelman1998}自由主義者要求這個問題必須可解，他的解法須指出現行根本法安排具有某種屬性，足以在道德上證成任何人對這個始終受爭議的武力體制所作的支持，使受其拘束的人遵守法律，不只出於避免痛苦的武力，也出於自覺持有的「respect for」〔尊重〕（註 70 引 Habermas 447--448 頁）。政權符合真正正義之事實若為公眾所知，固然足以引發尊重，這恰是合理的詮釋多元所排除的可能。\footcite[422 \& n.70]{michelman1998}

\section{暴露於全面意見衝擊之候選屬性}
\label{sec:40_paradox-4}

Michelman 另設一種可能：政權的根本法安排持續而可信地暴露於民主政治的批判嚴格檢驗之下，也許足以讓受其拘束的人出於尊重\ix{terms}{04@\zhlidxcat{正當性}!出於尊重}而遵守。\footcite[423]{michelman1998}他先說明，我們被迫同意兩件事：根本法的終極正確標準不能交由民主決定，而且須授權某個或某些官方機關不時決定並解釋根本法的內容；另一方面，我們不能承認任何官員的判斷是國家根本法應如何要求的可靠指引。在這兩點之外，人們仍可能相信最值得尊重的判斷，來自成員不斷暴露於「the full blast of the sundry opinions」〔各種意見之全面衝擊〕\ix{terms}{04@\zhlidxcat{正當性}!各種意見之全面衝擊}的官方機構，那些意見由社會各成員依其不同的經驗、位置與反省，自由而無拘束地形成。\footcite[423]{michelman1998}

Michelman 把這種信念的根據說成「a tight mix of epistemic and dignitary considerations」〔認識考量與尊嚴考量的緊密混合〕。認識的一面在於官員判斷的可靠性存疑，而有合理基礎期待判斷會因暴露於民主的「critical interaction」〔批判性互動〕\ix{terms}{01@\zhlidxcat{民主理論}!批判性互動}而在某種程度上得到啟發，「critical interaction」是 Michelman 括註所稱 Post\ix{persons}{post@Post, Robert} 對這個過程的貼切稱呼，註 71 參見 Post《Constitutional Domains》142--148 頁。尊嚴的一面見於他的兩個反問：既然如此，為何不把尊重保留給對公民顯示尊重的公共決定安排，即把公民視為政治啟蒙上潛在的稱職而真誠的貢獻者的安排？為何不把尊重保留給那些官方的努力，它們始終朝向使自己可受公共審議與辯論影響的目標，而那審議與辯論充分而公平地接納每個人對正義的意見？\footcite[423 \& n.71]{michelman1998}

\section{回應之不完整與社會的民主}
\label{sec:40_paradox-5}

這項回應，依 Michelman 自己的說法，對要求為支持現行根本法安排提出理由的人「can never be complete」〔永遠不可能完整〕。\footcite[423]{michelman1998}原因在於回應所稱「fully and fairly receptive」〔充分而公平地接納〕每個人意見的過程，必然是一個由法律構成的過程，而構成它的不只有關於政治代表與選舉的法律，還有關於民間結社、家庭、職場、言論、財產與媒體等的法律。因此，要以人人都能出於尊重而遵守為理由，證成對一套根本法安排的支持，就須確立這些法律合乎其在民主上應有的樣子，即構成一個大致「fair」〔公平〕或「undistorted」〔未受扭曲〕的民主政治溝通過程，範圍不限於立法與裁判的正式舞台，而及於整個公民社會。Michelman 以「social」〔社會的〕稱這種民主。\footcite[423--424]{michelman1998}

但依合理的詮釋多元的自由主義前提，這些法律是否符合「social-democracy-constitutive standard」〔構成社會的民主之標準〕，始終是有爭議而合理的歧見所在。Michelman 說：「We seem to have landed right back in the fix from which we seek escape」〔我們似乎又落回了我們想逃離的困境〕。\footcite[424]{michelman1998}

\section{民主地追求民主正義}
\label{sec:40_paradox-6}

424 頁接著說，所落之處未必完全相同。Michelman 期待自己已把政治正當性的問題重述為這樣：可以設想一些「conditions of apparent near-democracy, political and social」〔政治上與社會上顯然近乎民主的條件〕，在其中合理的公民得以互相要求同意，這個國家正善意而相當妥善地追求一項自治的政治計畫，該計畫「there can be no final end」〔沒有最終的終點〕，卻有「a foundational standard」〔一項基礎標準〕。\footcite[424]{michelman1998}

此一計畫即「The democratic pursuit of democratic justice」〔民主地追求民主正義〕\ix{terms}{01@\zhlidxcat{民主理論}!民主地追求民主正義}，其「first and constant requirement」〔首要而恆常的要求〕如下：對於基本政治安排是否正確所必須作成的決定，凡無法以合意達成者，仍須由這樣的機構作成，該機構始終有效地受一種形成中的公眾意見（public-opinion-in-formation）之壓力，這種公眾意見致力於使自身及其形成的法律與社會條件民主化。\footcite[424]{michelman1998}依這項要求，根本法的內容仍由官方機關決定而不由民主程序決定，原文第 V 部的結論因此不變；正當性的條件改以決定機關所受的壓力界定，即公眾意見的形成及其自我民主化。下一部論 Brennan，Michelman 在結尾說他提供了一項公平的測試（第~\ref{ch:50_brennan} 部）。

\section{課前可提的問題}
\label{sec:40_paradox-7}

\begin{enumerate}
\tightlist
\item 419 頁的無窮後退論證針對的是以民主的歷史作為尊重憲法的理由。以現行憲法已歷經多次修正為由而遵守，是否同樣須以強制施行該憲法為正當作為前提？
\item 421 頁所稱支持政權在道德上能否證成的問題，依 Fallon 的三種正當性概念（W04）屬於哪一種？423 頁以「epistemic」與「dignitary」兩項考量支撐的候選屬性，是否也涉及社會學意義的正當性？
\end{enumerate}
